\documentclass[12pt,hyperref=unicode]{beamer}

\usepackage[utf8]{inputenc}
\usepackage[turkish]{babel}
\usepackage[T1]{fontenc}

\usepackage{amsmath}
\usepackage{amssymb}
\usepackage{amsthm}
\usepackage{enumerate}
\usepackage{tikz}
\usepackage{transparent}
\usepackage{xcolor}
\usepackage{listings}
\usepackage{verbatim}
\usepackage{graphicx}

\usepackage{pgfplots}

\usetheme{Warsaw}
\usecolortheme{default}
\definecolor{nigdeyesili_acik}{RGB}{3, 150, 166}
\definecolor{nigdeyesili_koyu}{RGB}{143, 209, 217}
\setbeamercolor{structure}{fg=nigdeyesili_acik}

\author[Dr. Mehmet CANEVİ]{Arş.~Gör.~Dr.~M.~Canevi\inst{1}}

\institute{
    \inst{1}%
    Bilgisayar Mühendisliği\\
    Mühendislik Fakültesi
}
    
\date[2025] {Ders Notları, Ocak 2025}
\setbeamertemplate{footline}[frame number]

\logo{\transparent{0.4}\includegraphics[width = 20mm]{logo}}

\definecolor{mGreen}{rgb}{0,0.6,0}
\definecolor{mGray}{rgb}{0.5,0.5,0.5}
\definecolor{mPurple}{rgb}{0.58,0,0.82}
\definecolor{backgroundColour}{rgb}{0.95,0.95,0.92}

\lstloadlanguages{C}
\lstdefinestyle{CStyle}{
    backgroundcolor=\color{backgroundColour},   
    commentstyle=\color{mGreen},
    keywordstyle=\color{red},
    numberstyle=\tiny\color{mGray},
    stringstyle=\color{mGreen},
    basicstyle=\footnotesize,
    breakatwhitespace=false,         
    breaklines=true,                 
    captionpos=b,                    
    keepspaces=true,                 
    numbers=left,                    
    numbersep=2pt,                  
    showspaces=false,                
    showstringspaces=false,
    showtabs=false,                  
    tabsize=1,
    language=C
}
\lstset{style=CStyle}

\AtBeginDocument{\shorthandoff{=}}
\title[Ders 9] {P-tipi kontrolör I (Kuram)}
\begin{document}
%%%%%%%%%%%%%%%%%%%%%%%%%%%%%%%%%%%%%%%%%%%%%%%%%%%%%%%%%%%%%%%%%%%%%%%%%%%%%%%%
\frame{\titlepage}
\begin{frame}[fragile]{İçidekiler}
    \tableofcontents
\end{frame}
%%%%%%%%%%%%%%%%%%%%%%%%%%%%%%%%%%%%%%%%%%%%%%%%%%%%%%%%%%%%%%%%%%%%%%%%%%%%%%%%
\section{P-tipi kontrolör}
\begin{frame}[fragile]{Sistem}
    Birinci dereceden bir diferansiyel denklem
    \begin{equation}
        \dot{y}(t)+y(t)=u(t)
    \end{equation}
    olarak verilmiştir. $\mathcal{L}\{\dot{y}(t)\}=sY(s)$ kullanılarak
    \begin{equation}
    \begin{split}
        sY(s)+Y(s)&=U(s)\\
        (s+1)Y(s)&=U(s)\\
        \frac{Y(s)}{U(s)}&=\frac{1}{s+1}
    \end{split}
    \end{equation}
    elde edilir.
\end{frame}
%%%%%%%%%%%%%%%%%%%%%%%%%%%%%%%%%%%%%%%%%%%%%%%%%%%%%%%%%%%%%%%%%%%%%%%%%%%%%%%%
\section{P-tipi kontrolör}
\begin{frame}[fragile]{Sistem}
    Girişe $u(t)=1,\,t>0$($U(s)=\frac{1}{s}$) uygulanırsa
    \begin{equation}
    \begin{split}
        Y(s)&=\frac{1}{s(s+1)}\\
        &=\frac{1}{s}-\frac{1}{s+1}\\
        y(t)&=\mathcal{L}^{-1}\{\frac{1}{s}-\frac{1}{s+1}\}\\
        y(t)&=1-e^{-t}
    \end{split}
    \end{equation}
\end{frame}
%%%%%%%%%%%%%%%%%%%%%%%%%%%%%%%%%%%%%%%%%%%%%%%%%%%%%%%%%%%%%%%%%%%%%%%%%%%%%%%%
\begin{frame}[fragile]{Sistem yanıtı}
    \begin{figure}[!h]
        \centering
        \begin{tikzpicture}
        \begin{axis}[grid=both,xlabel={Zaman$(s)$},ylabel={$y(t)$},title={Basamak yanıtı},xmin=0,xmax=5,ymin=0,ymax=1.1,minor tick num=2]
            \addplot[black,scale=1,line width=1.2pt, domain=0:5, smooth]{1};
            \addplot[blue,scale=1,line width=1.2pt, domain=0:5, smooth]{1-exp(-x)};
        \end{axis}
        \end{tikzpicture}
    \end{figure}
    olarak elde edilir.
\end{frame}
%%%%%%%%%%%%%%%%%%%%%%%%%%%%%%%%%%%%%%%%%%%%%%%%%%%%%%%%%%%%%%%%%%%%%%%%%%%%%%%%
\end{document}